\begin{resumo}

    O sistema de gestão do PMOC de um órgão público estadual programa a manutenção preventiva de mais de mil aparelhos de ar-condicionado sem considerar os diferentes tempos de atendimento nem os deslocamentos entre prédios. A partir de tempos medidos em campo, o problema foi formulado como um sequenciamento de equipes paralelas, com tempos de preparação dependentes da sequência, jornada diária e limite de atendimentos por prédio. Para resolvê-lo, foi desenvolvido um modelo de programação inteira mista, utilizado na validação em instâncias pequenas, e implementadas quatro meta-heurísticas (IG, VNS, ILS e AG), cujos parâmetros foram calibrados por busca em grade e Optuna. Nos experimentos com 18 instâncias de 1000 a 2000 aparelhos, o VNS foi o único algoritmo que apresentou melhorias consistentes em relação à solução construtiva LPT+NEH. Essa melhoria foi estatisticamente significativa em comparação ao IG e ao AG, mas não ao ILS. No cenário de referência, a programação atual estabelece um prazo de 45 dias, enquanto a execução estimada pela regra vigente requer 74 dias. A programação obtida com os métodos propostos reduz esse período para 69 dias.

    \noindent\textbf{Palavras-chave}: PMOC. Manutenção preventiva. Máquinas paralelas. Meta-heurísticas. Programação inteira mista.

\end{resumo}

% --- resumo em inglês ---
\begin{resumo}[Abstract]
\begin{otherlanguage*}{english}

    The PMOC management system of a state public agency schedules preventive maintenance for more than one thousand air-conditioning units without accounting for differences in service times or travel between buildings. Based on times measured in the field, the problem was formulated as a parallel-team scheduling problem with sequence-dependent setup times, daily working hours and building-level limits on the number of services. A mixed-integer programming model was developed to validate the approach on small instances, and four metaheuristics (IG, VNS, ILS and GA) were implemented, with their parameters tuned through grid search and Optuna. In experiments involving 18 instances with 1,000 to 2,000 units, VNS was the only algorithm to consistently improve upon the LPT+NEH constructive solution. This improvement was statistically significant compared with IG and GA, but not with ILS. In the reference scenario, the current schedule sets a 45-day deadline, whereas execution under the existing scheduling rule is estimated to require 74 days. The proposed scheduling approach reduces this period to 69 days.


    \noindent\textbf{Keywords}: preventive maintenance. Parallel machines. Sequence-dependent setup times. Metaheuristics. Mixed-integer programming.
\end{otherlanguage*}
\end{resumo}
